\section{Discussion}
\label{sec:discussion}

\paragraph{Drought as financial risk, interconnection as reliability risk.} The clearest result is a separation
of risks. With the interconnector available, even a 40\% inflow reduction is met entirely by imports. The cost
of drought is financial: operating cost rises by 284~MEUR/yr between the normal and the severe-drought year.
Unserved energy appears only when transfer availability falls to about half, and it is energy-limited.
Short-duration storage can therefore reduce it only partially (13\% at 50\% availability). At an assumed VOLL of
5{,}000~EUR/MWh, however, even that partial reduction is worth more than the battery costs. Batteries thus
earn their keep here as reliability assets in joint stress events, not as arbitrageurs. This agrees with the concern of van Vliet et
al.~\cite{vanvliet2016vulnerability} that climate impacts on water-dependent generation are partly
compensated by trade and partly transmitted to neighbours. It also suggests that, for Albania, the joint
probability of domestic drought and reduced regional export capacity deserves more attention than domestic
drought alone. Because Balkan hydro systems share weather, this joint risk is plausible. Our scenario
probabilities for it are assumed, not estimated.

\paragraph{Solar as a no-regret option, conditional on export.} Under the base assumptions, PV is worth about
1.7$\times$ its annualised cost at the margin. It dominates all other investments at every budget and in every
scenario, which is why VSS is zero. Two mechanisms drive this. First, in a reservoir system PV saves water,
which the reservoir then moves to high-price hours, so PV inherits part of the reservoir's flexibility value.
Second, surplus PV can be exported at 85\% of the import price. The decline of renewable market value with
penetration described by Hirth~\cite{hirth2013market} is therefore weak here, but only because the export
route is assumed to absorb the surplus. When exports are restricted (Experiment~I2), this conclusion changes.
That restriction is the condition under which planning for uncertainty starts to matter.

\paragraph{Storage: timing, not volume.} Batteries reduce cost by shifting energy in time. In doing so they
\emph{increase} gross imports, by importing cheaply and serving the evening peak. Policy goals stated as
``reducing import dependence'' and ``reducing system cost'' can therefore point in different directions for
storage. Arbitrage alone does not recover battery costs here: savings are 13--27\% of annualised cost across the tested sizes and years. The
economic case for batteries rests instead on avoided shortage in joint drought and import-stress events
(Section~\ref{sec:sensitivity}), and so on how likely those events are and what VOLL is assumed.
The marginal analysis shows, however, that discharge power is worth more than its cost once some energy
capacity exists. This points to shorter-duration systems and to the reserve and ancillary services that the
model omits \cite{lund2015flexibility,kondziella2016flexibility}. The reservoir already supplies long-duration
flexibility: additional reservoir \emph{volume} has zero marginal value in every case, because annual inflow,
not storage size, is the binding resource.

\paragraph{Methodological lessons.} Two modelling choices changed answers by more than many of the parameter
uncertainties. First, temporal resolution: at 6-hourly steps battery value nearly vanishes and stochastic
planning had nothing to choose between (Section~\ref{sec:sensitivity}), consistent with
Poncelet et al.~\cite{poncelet2016temporal}. Second, how the boundaries of representative periods are handled
for a seasonal store. Pinning each week's terminal volume created spurious scarcity. That made a battery look
valuable for tail risk and reversed the ranking of two plans compared with full chronology. Valuing the end
volume at the coarse annual run's water value, a single-cut future-cost function as in SDDP
\cite{pereira1991sddp}, removed the artefact. This is a concrete example of the seasonal-storage problem in
representative-period models described by Kotzur et al.~\cite{kotzur2018seasonal}. Our recommendation is that
studies of hydro-dominated systems validate their period selection against full chronology, as we do here.

\paragraph{Model implications versus policy.} These are implications of a single-node model with synthetic
demand and prices. They are not recommendations for Albanian policy. A credible policy analysis would need
observed hourly load (ENTSO-E), regional price and transfer-capacity data, calibrated hydrology and a network
representation. The open pipeline is designed so that each of these can be substituted without code changes
to the optimisation core.
